\documentclass[12pt]{article}
\usepackage{amsmath,amssymb,amsfonts,amsthm}
\usepackage[margin=1in]{geometry}

\begin{document}

\begin{center}
{\Large\bfseries Chapter 4, Problem 15}\\[6pt]
Byron \& Fuller, \textit{Mathematics of Classical and Quantum Physics}
\end{center}

\bigskip

\textbf{Problem.} Show that the general Lorentz transformation worked out in Problem 14 may be expressed in vector notation as follows:
\[
\mathbf{r}' = \mathbf{r} + \frac{\gamma - 1}{v^2}(\mathbf{v}\cdot\mathbf{r})\mathbf{v} - \gamma\mathbf{v}t, \quad t' = \gamma\!\left(t - \frac{\mathbf{v}\cdot\mathbf{r}}{c^2}\right),
\]
where $\mathbf{v} = v\hat{n}$, $\beta = v/c$, and $\gamma = 1/\sqrt{1-v^2/c^2}$ is the relative velocity of the two frames. These formulas are a compact, convenient way to write the Lorentz transformation between systems with parallel axes moving with arbitrarily directed uniform relative velocity.

\bigskip
\textbf{Solution.}

From Problem 14, the general Lorentz boost with velocity $\mathbf{v} = v\hat{n}$ has components:
\[
\Lambda^{0}{}_{0} = \gamma, \quad \Lambda^{0}{}_{i} = -\beta\gamma\, n_i, \quad \Lambda^{i}{}_{0} = -\beta\gamma\, n_i, \quad \Lambda^{i}{}_{j} = \delta_{ij} + (\gamma-1)n_i n_j.
\]

\textbf{Time transformation:}
\[
ct' = \Lambda^{0}{}_{0}\,(ct) + \Lambda^{0}{}_{i}\,x^i = \gamma\, ct - \beta\gamma\, n_i x^i = \gamma\, ct - \gamma\frac{v}{c}\,(\hat{n}\cdot\mathbf{r}).
\]
Dividing by $c$:
\[
t' = \gamma\!\left(t - \frac{v(\hat{n}\cdot\mathbf{r})}{c^2}\right) = \gamma\!\left(t - \frac{\mathbf{v}\cdot\mathbf{r}}{c^2}\right). \checkmark
\]

\textbf{Spatial transformation:}
\[
x'^i = \Lambda^{i}{}_{0}\,(ct) + \Lambda^{i}{}_{j}\,x^j = -\beta\gamma\, n_i\,(ct) + \delta_{ij}x^j + (\gamma-1)n_i n_j x^j.
\]
\[
= x^i + (\gamma-1)(\hat{n}\cdot\mathbf{r})\,n_i - \gamma v\, n_i\, t.
\]

In vector form:
\[
\mathbf{r}' = \mathbf{r} + (\gamma-1)(\hat{n}\cdot\mathbf{r})\hat{n} - \gamma\mathbf{v}\,t.
\]

Since $\hat{n} = \mathbf{v}/v$, we have $(\hat{n}\cdot\mathbf{r})\hat{n} = \frac{(\mathbf{v}\cdot\mathbf{r})}{v^2}\mathbf{v}$. Therefore:
\[
\mathbf{r}' = \mathbf{r} + \frac{\gamma-1}{v^2}(\mathbf{v}\cdot\mathbf{r})\mathbf{v} - \gamma\mathbf{v}\,t. \checkmark
\]

This confirms the stated vector form of the general Lorentz transformation. \qed

\end{document}
